\section{2009级工科数学分析下 B卷}

\subsection{资料来源}
\begin{itemize}
  \item 试题：2009级计算机工科数学分析下 期末B.doc，DOC 正文已抽取；公式对象已按 Word 公式图片逐项恢复。
  \item 公式图片审计：\texttt{tmp/past-exam-equations/2009-B/manifest.csv}。
\end{itemize}

\subsection{试题}
诚信应考，考试作弊将带来严重后果！

华南理工大学 2009--2010 第二学期期末考试

《工科数学分析》（下）试卷（B）

注意事项：
\begin{enumerate}
  \item 考前请将密封线内填写清楚；
  \item 所有答案请直接答在试卷上；
  \item 考试形式：闭卷；
  \item 本试卷共五大题，满分 100 分，考试时间 120 分钟。
\end{enumerate}

\subsubsection*{一、填空题（共 5 小题，每小题 4 分，共 20 分）}
\begin{enumerate}
  \item 方程 $y''+5y'+4=0$ 的通解为 \underline{\hspace{4cm}}。
  \item 原点到平面 $x+\dfrac{y}{-4}+\dfrac{z}{3}=1$ 的距离为 \underline{\hspace{3cm}}。
  \item 函数 $z=\arctan\dfrac{y}{x}+\ln\sqrt{x^2+y^2}$ 在点 $(1,0)$ 处的梯度 $\operatorname{grad}z\big|_{(1,0)}=$ \underline{\hspace{3cm}}。
  \item 设 $L$ 是 $xOy$ 平面上任意一条分段光滑的闭曲线，则第二类曲线积分
  \[
  \oint_L(2xy+\cos x)\,dx+(x^2+\sin y)\,dy=
  \]
  \underline{\hspace{3cm}}。
  \item 设 $f(x)$ 是以 $2\pi$ 为周期的函数，它在 $[-\pi,\pi]$ 上的表达式是
  \[
  f(x)=
  \begin{cases}
  x+1,&-\pi\le x<0,\\
  x^2,&0\le x<\pi.
  \end{cases}
  \]
  若它的傅里叶级数的和函数为 $S(x)$，则 $S(0)=$ \underline{\hspace{3cm}}。
\end{enumerate}

\subsubsection*{二、选择题（共 5 小题，每小题 4 分，共 20 分）}
\begin{enumerate}
  \item 在 $R^3$ 中，方程 $x^2=4y$ 的图形是（\quad）。

  A. 抛物线；\quad B. 抛物柱面；\quad C. 椭圆抛物面；\quad D. 旋转抛物面。

  \item 曲面 $3x^2+y^2+z^2=12$ 上点 $M(-1,0,3)$ 处的切平面与平面 $z=0$ 的夹角是（\quad）。

  A. $\dfrac{\pi}{6}$；\quad B. $\dfrac{\pi}{4}$；\quad C. $\dfrac{\pi}{3}$；\quad D. $\dfrac{\pi}{2}$。

  \item 设 $f(x,y)$ 在点 $(x_0,y_0)$ 处的两个偏导数都存在，则（\quad）。

  A. $f(x,y)$ 在点 $(x_0,y_0)$ 必连续；\quad
  B. $f(x,y)$ 在点 $(x_0,y_0)$ 必可微；\quad
  C. $\displaystyle\lim_{x\to x_0}f(x,y_0)$ 与 $\displaystyle\lim_{y\to y_0}f(x_0,y)$ 都存在；\quad
  D. $\displaystyle\lim_{\substack{x\to x_0\\y\to y_0}}f(x,y)$ 存在。

  \item 设平面曲线 $L$ 为下半圆周 $y=-\sqrt{1-x^2}$，则曲线积分 $\displaystyle\int_L(x^2+y^2)\,ds$（\quad）。

  A. $0$；\quad B. $-1$；\quad C. $2$；\quad D. $\pi$。

  \item 幂级数 $\displaystyle\sum_{n=1}^{\infty}\frac{\ln(n+1)}{n}x^n$ 的收敛域为（\quad）。

  A. $(-1,1)$；\quad B. $[-1,1]$；\quad C. $[-1,1)$；\quad D. $(-1,1]$。
\end{enumerate}

\subsubsection*{三、计算题（本大题共 4 小题，每小题 8 分，共 32 分）}
\begin{enumerate}
  \item 求微分方程 $y''+y'+2y=0$ 的通解。
  \item 计算二重积分 $\displaystyle I=\iint_D\frac{x}{x^2+y^2}\,dxdy$，其中 $D$ 由曲线 $y=\dfrac{x^2}{2}$ 与直线 $y=x$ 围成。
  \item 计算曲线积分
  \[
  I=\int_L (x+y)^2\,dx-(x^2+y^2\sin y)\,dy,
  \]
  其中 $L$ 是抛物线 $y=x^2$ 上从点 $(-1,1)$ 到点 $(1,1)$ 的那一段。
  \item 求幂级数 $\displaystyle\sum_{n=0}^{\infty}\frac{n}{n+1}x^n$ 的收敛域及和函数。
\end{enumerate}

\subsubsection*{四、应用题（本大题共 2 小题，每小题 7 分，共 14 分）}
\begin{enumerate}
  \item 欲造一个无盖的长方体容器，已知底部造价为 $3$ 元/平方米，侧面造价为 $1$ 元/平方米。现想用 36 元造一个容积最大的容器，问应该怎样设计尺寸？
  \item 设 $\mathbf F(x,y)=xy^2\mathbf i+yf(x)\mathbf j$ 为保守场，其中 $f(x)$ 具有连续导数，且 $f(0)=0$。求曲线积分
  \[
  I=\int_{(0,0)}^{(1,1)}\mathbf F\cdot d\mathbf s.
  \]
\end{enumerate}

\subsubsection*{五、证明题（本大题共 2 小题，每小题 7 分，共 14 分）}
\begin{enumerate}
  \item 证明：圆柱螺线 $x=a\cos t,\ y=a\sin t,\ z=bt$ 的切线与 $z$ 轴的夹角为常数。
  \item 若正项级数 $\displaystyle\sum_{n=1}^{\infty}a_n$ 收敛，证明正项级数
  \[
  \sum_{n=1}^{\infty}\frac{\sqrt{a_n}}{n}
  \]
  也收敛。
\end{enumerate}

\subsection{答案与解析}

\subsubsection*{一、填空题}
\begin{enumerate}
  \item 令 $p=y'$，则 $p'+5p=-4$。解得
  \[
  p=C_1e^{-5x}-\frac45,
  \]
  故通解为
  \[
  y=C_2+C_1e^{-5x}-\frac45x.
  \]
  \item 平面化为 $x-\dfrac14y+\dfrac13z-1=0$，原点到平面的距离
  \[
  d=\frac1{\sqrt{1+\frac1{16}+\frac1{9}}}=\frac{12}{13}.
  \]
  \item 与 A 卷同题，
  \[
  \operatorname{grad}z\big|_{(1,0)}=(1,1).
  \]
  \item 由格林公式，
  \[
  \frac{\partial}{\partial x}(x^2+\sin y)-\frac{\partial}{\partial y}(2xy+\cos x)=2x-2x=0.
  \]
  对任意闭曲线积分为 $0$。
  \item $x=0$ 是跳跃点，故
  \[
  S(0)=\frac{f(0-)+f(0+)}2=\frac{1+0}{2}=\frac12.
  \]
\end{enumerate}

\subsubsection*{二、选择题}
\begin{enumerate}
  \item 方程与 $z$ 无关，是沿 $z$ 方向延伸的抛物柱面，选 B。
  \item 曲面法向量为 $\nabla(3x^2+y^2+z^2)=(-6,0,6)$，平面 $z=0$ 的法向量为 $(0,0,1)$。两平面夹角满足
  \[
  \cos\theta=\frac{|(-6,0,6)\cdot(0,0,1)|}{6\sqrt2}=\frac1{\sqrt2},
  \]
  故 $\theta=\dfrac{\pi}{4}$，选 B。
  \item 偏导数存在只保证沿坐标方向的一元极限存在。选 C。
  \item 下半圆周是单位圆的一半，且 $x^2+y^2=1$，故
  \[
  \int_L(x^2+y^2)\,ds=\int_L1\,ds=\pi.
  \]
  选 D。
  \item 半径为 $1$。$x=1$ 时为 $\sum\ln(n+1)/n$ 发散；$x=-1$ 时为交错级数 $\sum(-1)^n\ln(n+1)/n$，因 $\ln(n+1)/n\downarrow0$ 而收敛。故收敛域为 $[-1,1)$，选 C。
\end{enumerate}

\subsubsection*{三、计算题}
\begin{enumerate}
  \item 特征方程
  \[
  r^2+r+2=0,\qquad r=\frac{-1\pm i\sqrt7}{2}.
  \]
  故通解为
  \[
  y=e^{-x/2}\left(C_1\cos\frac{\sqrt7}{2}x+C_2\sin\frac{\sqrt7}{2}x\right).
  \]
  \item 区域为 $0\le x\le2,\ x^2/2\le y\le x$。于是
  \[
  I=\int_0^2\int_{x^2/2}^{x}\frac{x}{x^2+y^2}\,dy\,dx
  =\int_0^2\left[\arctan\frac{y}{x}\right]_{x^2/2}^{x}dx.
  \]
  因此
  \[
  I=\int_0^2\left(\frac{\pi}{4}-\arctan\frac{x}{2}\right)dx
  =\ln2.
  \]
  \item 与 A 卷同题，取 $x=t,\ y=t^2,\ -1\le t\le1$，得
  \[
  I=\int_{-1}^{1}(t^2+t^4)\,dt=\frac{16}{15}.
  \]
  \item 当 $|x|<1$，
  \[
  \sum_{n=0}^{\infty}\frac{n}{n+1}x^n
  =\sum_{n=0}^{\infty}x^n-\sum_{n=0}^{\infty}\frac{x^n}{n+1}
  =\frac1{1-x}+\frac{\ln(1-x)}{x}.
  \]
  端点 $x=1$ 时发散，$x=-1$ 时收敛，故收敛域为 $[-1,1)$。和函数为
  \[
  S(x)=\frac1{1-x}+\frac{\ln(1-x)}{x}\quad(-1<x<1),
  \]
  并有 $S(0)=0$，$S(-1)=\dfrac12-\ln2$。
\end{enumerate}

\subsubsection*{四、应用题}
\begin{enumerate}
  \item 设底面长宽为 $x,y$，高为 $h$。费用约束为
  \[
  3xy+2xh+2yh=36.
  \]
  最大化 $V=xyh$。由对称性取 $x=y=a$，则
  \[
  3a^2+4ah=36,\qquad V=a^2h=\frac{a(36-3a^2)}4.
  \]
  令 $V'=9-\dfrac94a^2=0$，得 $a=2$，进而 $h=3$。故底面为 $2\times2$，高为 $3$。
  \item 保守场条件为 $P_y=Q_x$，即
  \[
  2xy=yf'(x).
  \]
  故 $f'(x)=2x$，又 $f(0)=0$，得 $f(x)=x^2$。势函数满足
  \[
  \varphi_x=xy^2,\qquad \varphi_y=x^2y.
  \]
  取 $\varphi=\dfrac12x^2y^2$。于是
  \[
  I=\varphi(1,1)-\varphi(0,0)=\frac12.
  \]
\end{enumerate}

\subsubsection*{五、证明题}
\begin{enumerate}
  \item 螺线切向量为
  \[
  r'(t)=(-a\sin t,\ a\cos t,\ b).
  \]
  与 $z$ 轴方向 $k=(0,0,1)$ 的夹角 $\theta$ 满足
  \[
  \cos\theta=\frac{|r'(t)\cdot k|}{|r'(t)|}=\frac{|b|}{\sqrt{a^2+b^2}},
  \]
  为常数。
  \item 由不等式 $2\sqrt{uv}\le u+v$，取 $u=a_n$，$v=1/n^2$，得
  \[
  \frac{\sqrt{a_n}}{n}\le\frac12a_n+\frac1{2n^2}.
  \]
  因 $\sum a_n$ 与 $\sum1/n^2$ 均收敛，由比较判别法知
  \[
  \sum_{n=1}^{\infty}\frac{\sqrt{a_n}}{n}
  \]
  收敛。
\end{enumerate}
